% !TeX root = ../../Book2.tex
\chapter*{序言}
\addcontentsline{toc}{chapter}{序言}

一根弦的振动可以分成不同频率的振动，一团热量的扩散可以由一个核描述，而一个函数的局部粗糙程度也常能从它的频率分布中辨认。这些事实把 Fourier 分析与偏微分方程放在一起。但从一个漂亮的形式公式走到可靠的结论，中间还要回答许多问题：级数在哪种意义下收敛，微分与积分能否交换，边界条件如何进入算子，频率上的一个简单截断为何可能不再是有界操作。

\ProseParagraphBreak

本卷从这些具体问题出发，先建立 Fourier 变换的基本语言，再研究经典方程、弱解和调和分析中的估计。它承接第一卷的测度、积分、函数空间和 Sobolev 理论，但不要求读者事先熟悉偏微分方程或现代调和分析。读者应掌握多元微积分、线性代数、基本常微分方程，并能熟练使用 Lebesgue 收敛定理、\(L^p\) 空间和 Hilbert 空间的初步性质；涉及弱紧性、迹与嵌入时，正文会指出第一卷相应的位置。

\ProseParagraphBreak

第一编区分周期级数、\(L^1\) 变换、\(L^2\) 等距理论和分布变换，逐步建立反演与逼近。第二编先沿特征线解一阶方程，再通过分离变量、能量与基本解研究波动、热传导和势方程。这两编适合组成初次学习的主体：每遇到一个表示公式，都应检查初值、边界、收敛意义和唯一性，避免把形式求解与存在定理混为一谈。

\ProseParagraphBreak

第三编把经典方程转为能量空间中的问题。变分法解释方程从何而来，Lax--Milgram 方法给出弱存在性，正则性估计再回答弱解何时具有更多导数。这里最值得反复练习的是测试函数的选择：一个差商、一段截断或一个正部，会决定哪些项可以吸收，哪些假设真正进入估计。附录 H 进一步讨论低正则系数和非线性能量，但始终把局部有界、Harnack 比较与更高正则性分开说明。

\ProseParagraphBreak

第四、五编转向算子的尺度与振荡。极大函数、插值、奇异积分和平方函数提供基本估计；测不准、驻相、乘子、限制与 Bochner--Riesz 问题则说明局部化和曲率怎样影响频率分析。越到后面，读者越需要注意参数范围。一个端点、一个对偶指数或一个曲率条件的变化，可能使原来成立的估计失效。必要条件和反例因此与正面定理并列出现。

\ProseParagraphBreak

全卷内容超过一门课程的容量。侧重方程的读者可读前三编，并按需要查阅特殊函数、半群、色散估计与非线性附录；侧重调和分析的读者可在第一编之后接读第四、五编，回查所需的方程例子。Bessel、Hermite 与球谐函数的附录负责统一归一化和基础计算，其他附录分别补充 Hardy 空间、加权理论、集中紧性与微局部方法。各编的依赖图只画直接且重要的前置关系，不能代替具体定理中的假设。

\ProseParagraphBreak

本卷统一采用 \(e^{-2\pi i x\cdot\xi}\) 的 Fourier 变换约定。常数不会影响所有定性结论，却会影响基本解、球面核和演化相位，因此计算中仍须逐项保留。内积对第一变量线性，分布配对采用双线性约定；读别的书时，应先核对这几项，再比较公式。

\ProseParagraphBreak

每章习题附有解答，宜先独立完成，再检查自己的论证是否说明了合法性与边界情形。正文中的基础证明不依赖读者先做完某一道题。少数深层定理采用标明出处的证明概要，正文会说明省略的步骤和使用范围；它们适合在掌握主线之后继续追索，不能把一段概要读成完整证明。

\ProseParagraphBreak

参考资料各有用途。Fourier 基础可伴读 Stein 与 Shakarchi 的 \cite{Stein2003Fourier}，偏微分方程可结合 Evans 的 \cite{Evans2010Partial}，奇异积分及较高阶调和分析可沿 Stein 的两部专著 \cite{Stein1970Singular,Stein1993Harmonic} 继续。具体版本、参数和深层结论的出处放在相关正文处；这些阅读建议不表示几部书的论证范围与本卷完全相同。

\ProseParagraphBreak

读完一个证明后，试着把频率尺度、空间尺度和控制量单独写在纸上，再把函数平移或缩放一次。如果结论仍能正确变换，便往往已经理解了证明最重要的一部分。接下来再问：哪一个抵消、紧性或正性保证了估计？这会比只记住定理名字更有助于处理新的方程。

\bigskip
\begin{flushright}
R. EverStarine

2026年10月11日
\end{flushright}
